\BourbakiStarChapter{致读者}

\begin{enumerate}
\def\labelenumi{\arabic{enumi}.}
\item
  本系列丛书（卷目见第 ix 页和第 x 页）从数学的起点开始，并给出完整的证明。原则上，它不要求读者具备特定的数学知识，而只要求对数学推理有一定的熟悉程度，并具备一定的抽象思维能力。尽管如此，它尤其面向那些至少已掌握大学数学课程第一年或第二年内容的读者。
\item
  我们选择的阐述方法是公理化的和抽象的，并且通常由一般到特殊。这一选择是由本论著的主要目的所决定的，即为整个现代数学体系提供坚实的基础。为此，必须熟悉相当大量的非常一般的概念和原理。此外，证明的要求对主题内容强加了一个严格固定的顺序。因此，除非读者已经具有相当广泛的数学知识，否则某些考虑的用处不会立即显现；否则，他必须有耐心暂缓判断，直到时机成熟。
\item
  为了减轻这一不利之处，我们经常在正文中插入一些例子，这些例子涉及读者可能已经知道但尚未在本系列中讨论过的事实。这样的例子总是放在两个星号之间：*\ldots*。毫无疑问，大多数读者会发现这些例子有助于他们理解正文，并且即使在第一次阅读时也宁愿不略去它们。当然，从纯粹逻辑的观点来看，略去它们不会有任何不利之处。
\item
  本系列分为若干卷（此处称为“书”）。前六本书有编号，并且一般来说，正文中的每一陈述都只假定前面各卷中已经讨论过的结果为已知。这一规则在每本书内部都成立，但为了阐述的方便，这些书不再按连续顺序排列。在每本书（或这些书的每一章）的开头，读者将找到关于其与其他书逻辑关系的精确说明，从而能够确信不存在任何恶性循环。
\item
  每一章的逻辑框架由该章的定义、公理和定理组成。这些是主要需要记住以备后续使用的部分。较不重要的结果以及那些容易从定理推出的结果被标记为“命题”、“引理”、“推论”、“注记”等。那些在第一次阅读时可以略去的内容用小号字印刷。对特别重要的定理的评注偶尔以“附释”的名义出现。
\end{enumerate}

为了避免冗长的重复，有时引入仅在某一章或某一章中的某一节内有效的记号或缩写是方便的（例如，在一章只涉及交换环时，“环”一词总是指“交换环”）。这样的约定总是被明确提及，通常在其出现的章节的开头。

\begin{enumerate}
\def\labelenumi{\arabic{enumi}.}
\setcounter{enumi}{5}
\item
  正文中的某些段落旨在预先警告读者避免严重错误。这些段落在页边用符号 Z（“危险弯道”）标示。
\item
  练习的设计既使读者能够确信自己已消化了正文，又使他注意到那些在正文中没有位置但仍然有趣的结果。最难的练习带有符号 ¶。
\item
  总的来说，我们遵循了普遍接受的术语，除非有充分的理由偏离它。
\item
  我们特别努力始终使用严格正确的语言，而不牺牲简洁性。我们尽可能在正文中提请注意语言的滥用，没有这些滥用，任何数学文本都有变得迂腐、甚至不可读的风险。
\item
  由于原则上正文是对一个理论的教条式阐述，它通常不包含对文献的引用。参考文献汇集在历史注记中，通常位于每章末尾。这些注记在适当之处也包含对理论中未解决问题的指示。
\end{enumerate}

每个历史注记之后的参考书目通常只包含在所述理论的发展中具有最重要意义的书籍和原始论文。它绝不自诩完备；特别是，仅用于确定优先权问题的参考文献几乎总是被省略。

至于练习，我们一般认为没有必要指出它们的来源，因为它们取自许多不同的来源（原始论文、教科书、习题集）。

\par\noindent\textbf{11. 对本系列某一部分的引用如下给出：}\par

\begin{enumerate}
\def\labelenumi{\alph{enumi})}
\item
  如果引用的是同一节中提出的定理、公理或定义，则用其编号引用。
\item
  如果它们出现在同一章的另一节中，则在该引用中也引用该节。
\item
  如果它们出现在同一本书的另一章中，则引用该章和该节。
\item
  如果它们出现在另一本书中，则首先用其标题引用该书。
\end{enumerate}

结果概要用字母 R 引用；因此，集合论，R 表示“集合论的结果概要”。

